\documentclass[10pt,a4paper]{amsart}
\usepackage{amsmath,amssymb,amsthm,graphicx}
\usepackage{verbatim}
\newenvironment{reference}{\comment}{\endcomment}
%\newenvironment{reference}{}{}
\newenvironment{slogan}{\comment}{\endcomment}
\newenvironment{history}{\comment}{\endcomment}

% For commutative diagrams we use Xy-pic
\usepackage[all]{xy}

% We use 2cell for 2-commutative diagrams.
\xyoption{2cell}
% Theorem environments.
%
\theoremstyle{plain}
\newtheorem{theorem}[subsection]{Theorem}
\newtheorem{proposition}[subsection]{Proposition}
\newtheorem{lemma}[subsection]{Lemma}

\theoremstyle{definition}
\newtheorem{definition}[subsection]{Definition}
\newtheorem{example}[subsection]{Example}
\newtheorem{exercise}[subsection]{Exercise}
\newtheorem{situation}[subsection]{Situation}

\theoremstyle{remark}
\newtheorem{remark}[subsection]{Remark}
\newtheorem{remarks}[subsection]{Remarks}

\numberwithin{equation}{subsection}

% Macros
%
\def\lim{\mathop{\mathrm{lim}}\nolimits}
\def\colim{\mathop{\mathrm{colim}}\nolimits}
\def\Spec{\mathop{\mathrm{Spec}}}
\def\Hom{\mathop{\mathrm{Hom}}\nolimits}
\def\Ext{\mathop{\mathrm{Ext}}\nolimits}
\def\SheafHom{\mathop{\mathcal{H}\!\mathit{om}}\nolimits}
\def\SheafExt{\mathop{\mathcal{E}\!\mathit{xt}}\nolimits}
\def\Sch{\mathit{Sch}}
\def\Mor{\mathop{\mathrm{Mor}}\nolimits}
\def\Ob{\mathop{\mathrm{Ob}}\nolimits}
\def\Sh{\mathop{\mathit{Sh}}\nolimits}
\def\NL{\mathop{N\!L}\nolimits}
\def\CH{\mathop{\mathrm{CH}}\nolimits}
\def\proetale{{pro\text{-}\acute{e}tale}}
\def\etale{{\acute{e}tale}}
\def\QCoh{\mathit{QCoh}}
\def\Ker{\mathop{\mathrm{Ker}}}
\def\Im{\mathop{\mathrm{Im}}}
\def\Coker{\mathop{\mathrm{Coker}}}
\def\Coim{\mathop{\mathrm{Coim}}}

% Boxtimes
%
\DeclareMathSymbol{\boxtimes}{\mathbin}{AMSa}{"02}

%
% Macros for moduli stacks/spaces
%
\def\QCohstack{\mathcal{QC}\!\mathit{oh}}
\def\Cohstack{\mathcal{C}\!\mathit{oh}}
\def\Spacesstack{\mathcal{S}\!\mathit{paces}}
\def\Quotfunctor{\mathrm{Quot}}
\def\Hilbfunctor{\mathrm{Hilb}}
\def\Curvesstack{\mathcal{C}\!\mathit{urves}}
\def\Polarizedstack{\mathcal{P}\!\mathit{olarized}}
\def\Complexesstack{\mathcal{C}\!\mathit{omplexes}}
% \Pic is the operator that assigns to X its picard group, usage \Pic(X)
% \Picardstack_{X/B} denotes the Picard stack of X over B
% \Picardfunctor_{X/B} denotes the Picard functor of X over B
\def\Pic{\mathop{\mathrm{Pic}}\nolimits}
\def\Picardstack{\mathcal{P}\!\mathit{ic}}
\def\Picardfunctor{\mathrm{Pic}}
\def\Deformationcategory{\mathcal{D}\!\mathit{ef}}
\title{A formally smooth ring map which is not flat}
\author{The Stacks Project Authors}
\date{stacks-project:examples.tex:local-snapshot-20261012; GNU FDL 1.2 or later}
\begin{document}
\maketitle
\noindent\textit{Editorial note.} All delivered mathematical text is the author's own, copied from the pinned Stacks Project chapter file examples.tex; no statement or proof step was rewritten, reordered or removed.\par\medskip
\noindent\textit{Editorial note.} The counted proof is the author's complete argument: Companion of lemma-formally-etale-not-presented in the formally etale and smooth nonflat counterexample section.\par\medskip
\noindent\textit{Editorial note.} The label lemma-formally-smooth-nonflat names the counted statement, and the proof environment immediately following it in the source is delivered verbatim as its proof.\par\medskip
\noindent\textit{Editorial note.} The counted statement is carried by the author's own lemma environment in the chapter source.\par\medskip
\noindent\textit{Allowed standard prerequisites (editorial).}\begin{itemize}
\item Formally smooth morphism of rings: the lifting property against square-zero ideal extensions, namely every solid commutative diagram into R and R/I with I squared zero admits a dotted lift making the diagram commute
\item The group algebra k of the additive group of the rationals: the k-algebra with basis x-alpha indexed by rational alpha and multiplication x-alpha x-beta equals x of alpha plus beta, so x-zero is the unit
\item For a surjective ring map with kernel J, the obstruction to lifting against a square-zero extension A-prime to A with kernel I is the induced map J over J squared into I
\item A nonzerodivisor mapping to zero shows a ring map is not flat
\end{itemize}\medskip
\label{section-formally-smooth-nonflat}

\noindent
Let $k$ be a field. Consider the $k$-algebra $k[\mathbf{Q}]$.
This is the $k$-algebra with basis $x_\alpha, \alpha \in \mathbf{Q}$
and multiplication determined by $x_\alpha x_\beta = x_{\alpha + \beta}$.
(In particular $x_0 = 1$.)
Consider the $k$-algebra homomorphism
$$
k[\mathbf{Q}] \longrightarrow k, \quad
x_\alpha \longmapsto 1.
$$
It is surjective with kernel $J$ generated by the elements $x_\alpha - 1$.
Let us compute $J/J^2$. Note that multiplication by $x_\alpha$ on $J/J^2$
is the identity map. Denote $z_\alpha$ the class of $x_\alpha - 1$ modulo
$J^2$. These classes generate $J/J^2$. Since
$$
(x_\alpha - 1)(x_\beta - 1) = x_{\alpha + \beta} - x_\alpha - x_\beta + 1 =
(x_{\alpha + \beta} - 1) - (x_\alpha - 1) - (x_\beta - 1)
$$
we see that $z_{\alpha + \beta} = z_\alpha + z_\beta$ in $J/J^2$.
A general element of $J/J^2$ is of the form $\sum \lambda_\alpha z_\alpha$
with $\lambda_\alpha \in k$ (only finitely many nonzero). Note that if the
characteristic of $k$ is $p > 0$ then
$$
0 = pz_{\alpha/p} = z_{\alpha/p} + \ldots + z_{\alpha/p} = z_\alpha
$$
and we see that $J/J^2 = 0$. If the characteristic of $k$ is zero, then
$$
J/J^2 = \mathbf{Q} \otimes_{\mathbf{Z}} k \cong k
$$
(details omitted) is not zero.

\medskip\noindent
We claim that $k[\mathbf{Q}] \to k$ is a formally smooth ring map
if the characteristic of $k$ is positive. Namely, suppose given a
solid commutative diagram
$$
\xymatrix{
k \ar[r] \ar@{..>}[rd] & A \\
k[\mathbf{Q}] \ar[u] \ar[r]^\varphi & A' \ar[u]
}
$$
with $A' \to A$ a surjection whose kernel $I$ has square zero.
To show that $k[\mathbf{Q}] \to k$ is formally smooth we have to prove
that $\varphi$ factors through $k$. Since $\varphi(x_\alpha - 1)$ maps
to zero in $A$ we see that $\varphi$ induces a map
$\overline{\varphi} : J/J^2 \to I$ whose vanishing is the obstruction
to the desired factorization. Since $J/J^2 = 0$ if the characteristic
is $p > 0$ we get the result we want, i.e., $k[\mathbf{Q}] \to k$ is
formally smooth in this case. Finally, this ring map is not flat, for
example as the nonzerodivisor $x_2 - 1$ is mapped to zero.


\begin{lemma}
\label{lemma-formally-smooth-nonflat}
There exists a formally smooth ring map which is not flat.
\end{lemma}

\begin{proof}
See discussion above.
\end{proof}

\end{document}
